\begin{abstract}
As LLMs are increasingly deployed within \emph{agentic systems}, their capabilities depend not only on the model weights but also on the \emph{harness}: the prompts, tools, control flow, memory, and orchestration code surrounding them. This makes automated \emph{harness optimization} -- the iterative and evaluation-guided improvement of a harness by an AI system -- both an important route to improving AI systems and a demanding capability for AI systems themselves. Yet the community lacks a common protocol for measuring how well frontier LLMs perform at this task. We introduce \heb, a benchmark for end-to-end
harness optimization under expensive and stochastic evaluation. An
optimizer, an LLM paired with a coding harness, receives a
target agent's seed harness, graded evaluation feedback,
and a fixed target-evaluation budget. It edits the harness and nominates
a final candidate, which is scored by its normalized gain over the seed
on a held-out test partition that remains inaccessible throughout
search. A trusted execution environment enforces the evaluation
boundary, meters target-agent resource use, and preserves candidate
versions for audit. We evaluate 5 frontier LLMs as optimizers both under a shared coding harness
and under their native harnesses across
\NumPaperTasks\ downstream tasks, over \NumScoredCells scored runs. Experiment results show that optimizer
models separate more than the coding harnesses they act through, native
harnesses are not consistently superior, and gains vary substantially across tasks and seed regimes. These results establish harness optimization as a measurable and discriminative capability with large space for improvement.

\end{abstract}